%%  Title, author block (omitted under \ifblind), abstract and keywords.
%%  Shared by assembly with the title-page document (assemble_paper.py).
\begin{document}

\begin{frontmatter}

\title{Specification Surfaces for Incremental Forecast Performance:
Decomposition and Simultaneous Inference for Nested Probability Forecasts}

%%  ROUND TWENTY-SIX.  The Guide for Authors asks for a complete affiliation
%%  address and an ORCID, and the referee found neither.  The city and country
%%  are here; the ORCID is registered through the submission system, where
%%  Elsevier collects it from the corresponding author's profile rather than
%%  from the manuscript source.
\ifblind\else
\author[ind]{Sudhir Dixit\corref{cor1}}
\cortext[cor1]{Corresponding author.}
\ead{sudhir.dixit1@gmail.com}
\address[ind]{Independent Researcher, \affiliationStreet, \affiliationCity, \affiliationRegion, \affiliationCountry}
\fi

\begin{abstract}
Forecast evaluation asks whether a predictor improves out-of-sample
probability forecasts over a nested benchmark, answering with one score
difference dependent on unstated choices: the benchmark, model class,
split, forecast origin, predictor completeness and scoring rule. We treat
the improvement as a surface over declared choices, with three
reporting objects: an exact variance decomposition separating
forecaster choices from split-to-split variation; simultaneous bootstrap
bands, their family-wise coverage measured by simulation; and resolution
regions, where its sign is determined. Across \nPairs\ binary forecasting tasks from
\nLogs\ public event logs, mostly IT-service and
administrative, the evaluation split and its interactions carry a median
\shareResamplingStratumPct\ of the variance, and resolved specifications (\nPairsMadResolved\ tasks) lie a
median \madResolvedMedian\ in area under the ROC curve from the conventional
report. In a bank
service-desk case study, the predictor improves reassignment forecasts at
first touch, not at incident creation, once the benchmark holds what is
assumed recorded by then.
\end{abstract}

%%  IJF RETARGET.  The guide asks for at least five keywords, recommends
%%  its own list (forecasters.org/ijf/keywords) and says there is no point
%%  repeating words that are in the title: all eight are from the list
%%  (checked on 2026-09-25 against the web.archive.org copy of that page
%%  dated 2025-06-23; the live page sits behind a browser check).  "Event
%%  logs", the paper's own term, was the eighth until then; it is not an IIF
%%  keyword, it indexes the paper under business-process analytics rather
%%  than forecasting, and the abstract already names the data.  It was
%%  replaced by "Forecasting practice", because the object the paper
%%  criticises and replaces is a reporting practice (one score difference
%%  quoted for a predictor).  "Uncertainty" and
%%  "specification surface" repeat the title and were replaced by
%%  "Decomposition" (the functional-ANOVA decomposition) and "Statistical
%%  tests" (the simultaneous max-t bands and the planned contrasts under
%%  Holm).  2026-09-25: the title gained "Decomposition" on 2026-09-24, so
%%  that keyword became a repeat; it is replaced by "Error measures" (IIF
%%  list; the scoring rule is a design axis).  The abstract spells out
%%  the area under the ROC curve: the guide asks for no jargon there.  "predictive process monitoring" was removed in round
%%  twenty-seven and stays out: Section 11 says prefix length, bucketing and
%%  sequence encoding are not axes of this design space, and a keyword is a
%%  claim of topical membership.
\begin{keyword}
Evaluating forecasts \sep Error measures \sep Bootstrapping \sep
Probability forecasting \sep Robustness \sep Decision making \sep
Statistical tests \sep Forecasting practice
\end{keyword}

\end{frontmatter}
